% ----------------------------------------------------------
% Exemplo de capítulo sem numeração, mas presente no Sumário
\chapter*[Introdução]{Introdução}
\addcontentsline{toc}{chapter}{Introdução}
% ----------------------------------------------------------

O câncer de colo de útero é uma neoplasia maligna do epitélio de revestimento do órgão, considerado um grave problema de saúde pública global, especialmente prevalente em regiões de baixa e média renda. No Brasil,  é o terceiro tipo mais comum de câncer entre as mulheres, excluindo-se os casos de câncer de pele não melanoma. O carcinoma epidermoide é o subtipo mais incidente, correspondendo a 90\% dos casos. A infecção persistente por tipos oncogênicos do Papiloma Vírus Humano (HPV) é responsável por 99\% dos casos de câncer cervical (INCA,2021). Com isso, é uma doença amplamente prevenível graças a intervenções eficazes como a vacinação contra o HPV, e o rastreamento, com destaque para os métodos baseados na detecção do vírus.

Em resposta a essa realidade, a Organização Mundial da Saúde (OMS) estabeleceu em 2020 a Iniciativa Global para a Eliminação do Câncer Cervical. O objetivo é reduzir a incidência para menos de 4 casos por 100.000 mulheres-ano em todos os países. Para atingir essa meta até 2030, a estratégia da OMS foca em três pilares: vacinar 90\% das meninas até os 15 anos, rastrear 70\% das mulheres com um teste de alta performance pelo menos duas vezes até os 45 anos, e tratar 90\% das mulheres diagnosticadas com pré-câncer ou câncer cervical. Para o sucesso da iniciativa, é crucial que haja monitoramento contínuo e aprimorado para identificar lacunas e direcionar as ações (OMS,2020) (ARBYN, et al. 2020).

Apesar dos esforços, a incidência e a mortalidade do câncer de colo de útero ainda estão muito acima da meta da OMS em muitas regiões, especialmente naquelas com menores níveis de desenvolvimento humano, evidenciando marcantes desigualdades geográficas e socioeconômicas. Cidades como Manaus (Amazonas) ainda enfrentam obstáculos substanciais. Apesar da disponibilidade de exames citológicos gratuitos desde 1990 e imunização contra o HPV desde 2013, a capital amazonense mantém uma das maiores incidências de câncer cervical no mundo, refletindo limitações no rastreamento, cobertura vacinal e monitoramento das cepas virais circulantes (TORRES et al, 2021). 

Para mudar esse cenário, a implementação de programas de rastreamento organizados, aliados ao acesso a tratamento acessível e a uma maior cobertura vacinal contra o HPV, será fundamental para reduzir a incidência e a mortalidade nas próximas décadas. É essencial aplicar as novas ferramentas de prevenção e controle de forma criteriosa, considerando o contexto de cada local, e continuar monitorando a doença globalmente, especialmente entre as populações mais vulneráveis, pois os fatores socioeconômicos desempenham um papel crucial nas desigualdades do câncer (SINGH et al, 2023).

Avanços recentes na região incluem a aplicação do teste OncoE6™, voltado para a detecção da oncoproteína E6 do HPV-16 e HPV-18, em populações ribeirinhas e de difícil acesso no Amazonas. Tal abordagem demonstrou-se viável e altamente específica, evidenciando o potencial de inovações capazes de mitigar barreiras logísticas e, ao mesmo tempo, permitir uma cobertura de triagem expandida, um aspecto fundamental para mulheres que, diante de resultados anormais, enfrentam dificuldades para obter acesso à colposcopia e ao tratamento especializado (Torres et al., 2018).
Paralelamente, estudos genômicos locais identificaram mutações missense em genes provenientes de variantes do HPV-16 que circulam especificamente no estado do Amazonas. No gene E6, as mais frequentes observadas foram T7392G (L90V) — leucina por valina (45,5\%), C7377T (H85Y) — histidina por tirosina (39,4\%), e G7187T (Q21H) — glutamina por histidina (39,4\%). No gene E7, foram encontradas: G7810A (G43E) — glicina por glutamato, e A7689G (N29S) — asparagina por serina. (RONDON et al, 2024). Tais alterações podem modificar propriedades estruturais e funcionais dessas oncoproteínas, impactando diretamente sua interação com genes supressores de tumor, como p53 e p16INK4a.

A ausência de dados estruturais tridimensionais sobre essas variantes locais representa uma importante lacuna científica, com implicações no desenvolvimento de vacinas personalizadas, terapias-alvo e métodos diagnósticos mais precisos. Frente a esse cenário, o presente projeto propõe o uso de tecnologias emergentes de inteligência artificial (IA) - em especial, o software AlphaFold - para modelar, simular e prever interações estruturais das oncoproteínas E6 e E7 mutadas com alvos celulares de importância oncológica. Complementando essa abordagem, serão realizados docking molecular, predição de epítopos antigênicos e análise de interação com anticorpos para clonagem in silico de um potencial teste diagnóstico, além da disponibilização dos dados em uma base pública.
Dessa forma, a presente proposta visa investigar mutações regionais e mundiais do HPV-16 associadas ao câncer cervical, empregando tecnologias computacionais avançadas aplicadas à modelagem molecular e simulações biofísicas, com ênfase no desenvolvimento de soluções inovadoras para o diagnóstico e prevenção da doença, especialmente adaptados à realidade epidemiológica da região amazônica. 
Com forte potencial de inovação translacional, o projeto utiliza da bioinformática para transformar dados genéticos locais em aplicações biomédicas relevantes e compactuam com os objetivos da OMS para erradicação do câncer de colo cervical até 2030. 


\section{Via da beta-catenina}

A via Wnt desempenha papel central na regulação da proliferação e diferenciação celular, estando frequentemente desregulada no câncer colorretal (CCR) \cite{Liu2022}. A $\beta$‑catenina é o principal efetor dessa via, atuando como cofator transcricional ao interagir com proteínas como TCF4, BCL9 e LEF1². Essas interações promovem a ativação de genes relacionados à progressão tumoral e resistência terapêutica³. O bloqueio seletivo dessas interfaces
desponta como estratégia promissora para o desenvolvimento de novas terapias‑alvo no CCR, especialmente por meio de abordagens in silico, que permitem prever afinidades e potenciais pontos de intervenção \cite{Liu2022}. Este estudo comparou a estabilidade e a afinidade de ligação dos complexos $\beta$‑catenina/TCF4, $\beta$‑catenina/BCL9 e $\beta$‑catenina/LEF1, buscando identificar possíveis alvos terapêuticos na via Wnt.